Fix \(m,q,\iota,\delta\), the subclass \(\mathcal C\), and the constants in the proposed statement.  Constants below may depend on these fixed quantities and on \(L,A,r_*,\underline c,\overline c\), but not on \(n,P,t,x\).  The proof gives the upper rate, the valid contained-bump converse, the independent angular-family converse, and the common-density comparison.

\emph{Identification, rank, and support.}
For \(x\in\mathcal B_{P,\delta}\), put \(S_{P,x}=|D_{P,x}|\).  By \(\mathrm{lem:radial\text{-}coarea\text{-}identification}\), the signed-radius conditional mean is \(N_{t,P,x}(s)/Q_{t,P,x}(s)\) for almost every positive radius.  The assumed bound \(Q_{t,P,x}(s)\geq\underline c s\) makes every ratio in the proof well defined for \(0<s\leq r_*\).  Since each polynomial term and the remainder vanish at zero,
\[
\lim_{s\downarrow0}\frac{N_{t,P,x}(s)}{Q_{t,P,x}(s)}=\mu_{t,P}(x),
\]
and the same lemma yields the identifying formula for the separately defined causal target \(\tau_P(x)=\mu_{1,P}(x)-\mu_{0,P}(x)\).

For \(0<h\leq r_*\), polar coarea and the strict arm convention give
\[
\mathbb E_P\widehat G^{(m)}_{t,x}(h)
=\int_0^1K_\triangle(u)v_m(u)v_m(u)^\top\frac{Q_{t,P,x}(hu)}h\,du.
\]
Therefore, for every \(a\in\mathbb R^m\),
\[
\underline c\,a^\top G_ma
\leq a^\top\mathbb E_P\widehat G^{(m)}_{t,x}(h)a
\leq\overline c\,a^\top G_ma.
\]
If \(a\ne0\), the polynomial \(a^\top v_m(u)\) cannot vanish on an interval, while \(uK_\triangle(u)>0\) on \((0,1)\).  Hence \(G_m\) is positive definite and
\[
\lambda_{\min}\{\mathbb E_P\widehat G^{(m)}_{t,x}(h)\}
\geq\underline c\lambda_{\min}(G_m)=\lambda_m(\underline c)>0.
\]
This verifies the population inverse and cutoff conditions.  The upper radial bound also gives kernel-window probability at most \(\overline c h^2/2\).

Write
\[
g_{t,P,x}(s)=\frac{N_{t,P,x}(s)}{Q_{t,P,x}(s)}
=\mu_{t,P}(x)+\sum_{k=1}^{m-1}b_{t,k,P}(x)s^k+R_{t,P,x}(s),
\qquad |R_{t,P,x}(s)|\leq As^m,
\]
and define
\[
\beta_{t,P,x}(h)=
\bigl(\mu_{t,P}(x),hb_{t,1,P}(x),\ldots,h^{m-1}b_{t,m-1,P}(x)\bigr)^\top.
\]
For the observations in arm \(t\), algebra gives
\[
\widehat r^{(m)}_{t,x}(h)
=\widehat G^{(m)}_{t,x}(h)\beta_{t,P,x}(h)+B_{t,x}(h)+Z_{t,x}(h),
\]
where \(B_{t,x}(h)\) replaces \(Y_i\) by \(R_{t,P,x}(S_{P,x,i})\), and \(Z_{t,x}(h)\) replaces it by \(Y_i-g_{t,P,x}(S_{P,x,i})\).  The equality holds almost surely because the conditional-mean identity is needed only off a radial null set.

The centered radial error is uniformly sub-Gaussian.  Indeed,
\[
Y-g_{t,P,x}(S_{P,x})=\{Y-\mu_{t,P}(X)\}+\{\mu_{t,P}(X)-g_{t,P,x}(S_{P,x})\}.
\]
Conditional on \(X\), the first term has the declared sub-Gaussian mgf.  Conditional on signed radius, the second is centered and lies in \([-2L,2L]\), because the H\"older ball bounds \(\lVert\mu_{t,P}\rVert_\infty\) by \(L\).  Convexity of the exponential, applied to the chord joining the endpoints of that interval, gives the bounded-centered mgf bound.  Iterated conditioning consequently yields
\[
\mathbb E_P\!\left[e^{\lambda\{Y-g_{t,P,x}(S_{P,x})\}}\mid D_{P,x}\right]
\leq e^{C_L\lambda^2/2},
\qquad \lambda\in\mathbb R.
\]

\emph{Uniform concentration and bias.}
Let \(\varepsilon_n=n^{-3}\).  Compact support places the trimmed boundary inside \([-L,L]^2\), so an ambient grid supplies an \(\varepsilon_n\)-net with at most \(C_Ln^6\) points.  At a net point, Gram summands have operator norm at most \(C_mh^{-2}\) after the empirical averaging normalization and second-moment scale at most \(C_{m,\overline c}h^{-2}\); the preceding radial-error mgf gives the same variance scale for each coordinate of \(Z_{t,x}(h)\).  Applying exponential Markov optimization to the iid product mgfs, and then a union bound over two arms, the fixed matrix dimension, and the polynomial-size net, gives, for any fixed \(B>0\), an event \(E_n(P)\) satisfying
\[
\sup_{P\in\mathcal C}P\{E_n(P)^c\}\leq n^{-B}
\]
on which, at every net point,
\[
\left\|\widehat G^{(m)}_{t,x}(h)-\mathbb E_P\widehat G^{(m)}_{t,x}(h)\right\|_{\mathrm{op}}
\leq\frac14\lambda_m(\underline c),
\qquad
\left|e_1^\top\{\widehat G^{(m)}_{t,x}(h)\}^{-1}Z_{t,x}(h)\right|
\leq C\sqrt{\frac{\log n}{nh^2}}.
\]
The first display and the population bound make every inverse legitimate and put the estimator on its active branch.

For each fixed \(P\), the \(C^1\) image \(\mathcal B_P=\gamma_P([0,1])\) has finite length, hence two-dimensional Lebesgue measure zero.  Absolute continuity gives \(P(X_i\in\mathcal B_P)=0\).  Thus, almost surely, as \(x\) moves on the boundary the strict arm indicator remains fixed for each observation.  The maps \(u\mapsto K_\triangle(u)v_m(u)\) and \(u\mapsto K_\triangle(u)v_m(u)v_m(u)^\top\), on \([0,\infty)\) with their zero continuation beyond one, are Lipschitz.  Together with
\[
\bigl|\lVert X_i-x\rVert_2-\lVert X_i-x'\rVert_2\bigr|\leq\lVert x-x'\rVert_2,
\]
this bounds Gram oscillation within a mesh cell by \(C\varepsilon_nh^{-3}\) and raw outcome-vector oscillation by \(C\varepsilon_nh^{-3}n^{-1}\sum_i|Y_i|\).  Weyl's inequality, followed by \(A^{-1}-B^{-1}=A^{-1}(B-A)B^{-1}\), extends cutoff activity and the estimator bound to every boundary point.  The target oscillation obeys \(|\mu_{t,P}(x)-\mu_{t,P}(x')|\leq L\lVert x-x'\rVert_2\).  No uniform chart derivative, speed, reach, or length bound is used.

On the kernel support, \(|R_{t,P,x}(S)|\leq Ah^m\).  The first diagonal Gram entry bounds the associated empirical kernel mass, while the empirical inverse norm is at most \(2/\lambda_m(\underline c)\).  The polynomial-reproduction identity therefore gives on \(E_n(P)\)
\[
\sup_{x\in\mathcal B_{P,\delta}}
|\widehat\mu^{(m),\mathrm{LP}}_{t,h,\underline c}(x)-\mu_{t,P}(x)|
\leq C\left\{h^m+\sqrt{\frac{\log n}{nh^2}}\right\}.
\]
Subtracting arms gives the same order for the causal curve.

On a failed-cutoff branch the estimate is zero.  On an active branch, \(\mathrm{lem:stabilized\text{-}distance\text{-}local\text{-}polynomial\text{-}admissible}\) gives
\[
|\widehat\mu^{(m),\mathrm{LP}}_{t,h,\underline c}(x)|
\leq Ch^{-2}\overline Y_n,
\qquad \overline Y_n=n^{-1}\sum_i|Y_i|.
\]
The H\"older and conditional tail envelopes imply \(\sup_P\mathbb E_P\overline Y_n^2\leq C\).  Cauchy--Schwarz and the arbitrary choice of \(B\) make the complement contribution negligible.  Hence, for \(h\leq r_*\),
\[
\sup_{P\in\mathcal C}\mathbb E_P\sup_{x\in\mathcal B_{P,\delta}}
|\widehat\tau^{(m),\mathrm{LP}}_{n,h}(x)-\tau_P(x)|
\leq C\left\{h^m+\sqrt{\frac{\log n}{nh^2}}\right\}.
\]
Balancing the terms gives
\[
h_n\asymp\left(\frac{\log n}{n}\right)^{1/(2m+2)},
\qquad
h_n^m\asymp\left(\frac{\log n}{n}\right)^{m/(2m+2)}.
\]
For fixed \(m\), \(h_n\leq r_*\) eventually and \(nh_n^2/\log n\to\infty\), as required by concentration.  The admissibility lemma makes this one explicit member of \(\mathcal T_n^{\mathrm{dist}}\), proving the estimator and minimax upper bounds.

\emph{Contained-bump converse when \(m\geq q\).}
Use \([-1/2,1/2]^2\), uniform density, the horizontal interface, and the half-plane assignment.  Choose \(r_*>0\) smaller than the trimmed boundary's distance to the vertical sides.  Then, for either arm and every trimmed center,
\[
Q_{t,P,x}(s)=\pi s,
\qquad 0<s\leq r_*.
\]
Pack \(M\asymp h_n^{-1}\) centers \(x_1,\ldots,x_M\) at separation \(3h_n\).  For a fixed \(C_c^\infty\) bump \(\phi\) with \(\phi(0)=1\), set
\[
\Delta_n=\eta h_n^m,
\qquad
\mu_{0,j}=0,
\qquad
\mu_{1,j}(z)=\Delta_n\phi((z-x_j)/h_n),
\]
and use independent standard Gaussian potential-outcome errors, deterministic assignment, and consistency.  For \(0\leq k\leq q\),
\[
\lVert D^k\mu_{1,j}\rVert_\infty\leq C_\phi\eta h_n^{m-k}\leq C_\phi\eta
\]
because \(m\geq q\) and eventually \(h_n\leq1\).  Small fixed \(\eta\) therefore gives the \(q\)-H\"older envelope.  Compact support, density bounds, the smooth injective chart, and the sub-Gaussian condition are immediate, so the laws belong to \(\mathcal P_{\mathrm{coarea}}\).

The same derivative bound through order \(m\), followed by multivariate Taylor expansion at an arbitrary trimmed center and integration over each half-circle, produces uniformly bounded coefficients \(b_{t,k,P}(x)\) and remainder at most \(C\eta s^m\).  Thus the radial conditions are proved for this submodel, not hidden in its name.  The target separation at its active center is \(\Delta_n\).  Conditional Gaussian relative entropy and bump support area \(O(h_n^2)\) give
\[
\operatorname{KL}(P_j^{\otimes n},P_0^{\otimes n})
=\frac n2\mathbb E\!\left[\mu_{1,j}(X)^2\mathbf1\{T=1\}\right]
\leq Cn\Delta_n^2h_n^2=C\eta^2nh_n^{2m+2}\leq C'\eta^2\log n.
\]
Since \(\log M=\{(2m+2)^{-1}+o(1)\}\log n\), choose \(\eta\) so the last display is at most \(\alpha\log M\), \(\alpha<1/4\).  If \(J\) is uniform on the alternatives, the log-sum inequality yields \(I(J;O_{1:n})\leq\alpha\log M\).  For every decoder error event \(E\),
\[
\log M=I(J;O_{1:n})+H(J\mid O_{1:n})
\leq\alpha\log M+\log2+P(E)\log M,
\]
so \(P(E)\geq1-\alpha-o(1)\).  Sup-norm error below \(\Delta_n/3\) would identify \(J\) by maximizing the estimates over the separated centers.  Thus even the full-location experiment has risk at least \(c\Delta_n\), and the coordinatewise class cannot do better.  Since \(\Delta_n\asymp(\log n/n)^{m/(2m+2)}\), this matches the upper rate.  When \(m<q\), the same candidate has a derivative of order \(q\) of size \(h_n^{m-q}\to\infty\), so it is not a valid member of the declared H\"older class; no converse is claimed in that regime.

\emph{Angular-family converse.}
For any class containing the family in \(\mathrm{lem:angular\text{-}density\text{-}packing}\), Poissonization at mean \(2n\), the disjoint-cell decomposition, and the coordinatewise-overlap direct-product lemma give probability at least \(1/2-o(1)\) of some bit error because adjacent separation is \(c_1a_n\) and local relative entropy is at most \(2\alpha\log M\), with \(2\alpha<1\).  Midpoint thresholding and marked de-Poissonization therefore give the coordinatewise-distance lower frontier \(ca_n\).

\emph{One common density with two different joint remainders.}
Take an interior point \(x=0\) of a horizontal interface in a square and a radial continuous cutoff \(\eta_0\) equal to one near zero and supported in a disk strictly inside the square.  For \(0<\alpha<1\), define
\[
f(z)=c_0\left\{1+\eta_0(\lVert z\rVert_2)\lVert z\rVert_2^\alpha\frac{z_1}{\lVert z\rVert_2}\right\},
\]
with the fraction set to zero at \(z=0\).  The perturbation tends to zero at the origin, integrates to zero on each complete circle, and is supported away from the square boundary.  Taking the disk radius below one makes \(f\) bounded, positive, continuous, and normalized for the uniform base constant \(c_0\).  With a locally constant regression \(\mu\equiv c\), one has \(N(s)=cQ(s)\) exactly.  With the affine regression \(\mu(z)=z_1\), on the treated upper half-circle and for all sufficiently small \(s\),
\[
\frac{N(s)}{Q(s)}
=\frac{s\int_0^\pi\cos\theta\{1+s^\alpha\cos\theta\}\,d\theta}
{\int_0^\pi\{1+s^\alpha\cos\theta\}\,d\theta}
=\frac12s^{1+\alpha}.
\]
If an order-two remainder held, then \(|s^{1+\alpha}/2-b_1s|\leq As^2\).  Division by \(s\) and passage to zero first force \(b_1=0\); division by \(s^2\) then gives \(s^{\alpha-1}/2\leq A\), a contradiction.  Thus the same density supports an exact constant ratio for one smooth regression and failure of the order-two remainder for another.  This proves only the stated existential non-determination of the joint remainder and not a categorical impossibility theorem about every density-only threshold.